\section{Ingeniería de Requisitos}
\label{sec:requisitos}

Este capítulo resume el trabajo de ingeniería de requisitos documentado en detalle en
\texttt{docs/requisitos/} (SRS v1.0.0, \texttt{docs/requisitos/SRS-v1.0.0.pdf}), que constituye un
capítulo dedicado y separado del diseño arquitectónico (Sección~\ref{sec:arquitectura}), conforme al
criterio D0R de esta guía.

\subsection{Requisitos funcionales y no funcionales}

El sistema especifica 12 requisitos funcionales (Tabla~\ref{tab:rf-resumen}) y 4 requisitos no
funcionales (rendimiento, seguridad, usabilidad, mantenibilidad). Cada RF está documentado como una
historia de usuario en formato Connextra (``Como/Quiero/Para''), evaluada contra el criterio INVEST, con
criterios de aceptación expresados en Gherkin (\texttt{docs/requisitos/historias/}), y como un caso de
uso ``fully dressed'' con la notación de nivel de meta de Cockburn y un diagrama de secuencia
(\texttt{docs/requisitos/casos-de-uso/}).

\begin{table}[H]
\centering
\small
\begin{tabular}{p{1.2cm}p{5.5cm}p{1.6cm}p{2.5cm}}
\toprule
\textbf{RF} & \textbf{Título} & \textbf{Prioridad} & \textbf{Verificación} \\
\midrule
RF-01 & Autenticación segura & Must & Verificado (test real) \\
RF-02 & Registro de solicitud & Must & Verificado (test real) \\
RF-03 & Asignación de jurados & Must & Verificado (test real) \\
RF-04 & Programación de cronograma & Must & Verificado (test real) \\
RF-05 & Evaluación por rúbrica & Must & Verificado (test real) \\
RF-06 & Generación de actas & Must & No verificado \\
RF-07 & Firma digital de actas & Should & No verificado \\
RF-08 & Notificaciones & Could & No verificado \\
RF-09 & Reportes de defensas & Could & No verificado \\
RF-10 & Gestión de salas & Should & No verificado \\
RF-11 & Gestión de usuarios & Must & Verificado (test real) \\
RF-12 & Carga de anteproyecto & Must & Verificado (test real) \\
\bottomrule
\end{tabular}
\caption{The 12 functional requirements, their MoSCoW priority, and their real automated
verification status (source: \texttt{docs/trazabilidad/matriz.csv} v1.0.0).}
\label{tab:rf-resumen}
\end{table}

\subsection{Calidad de los requisitos: checklist INCOSE}

Se aplicaron las nueve características individuales de calidad de INCOSE \cite{incose2023} (necesario,
apropiado, no ambiguo, completo, singular, factible, verificable, correcto, conforme) a los 12 RF, y las
seis características de calidad del conjunto completo (\texttt{docs/checklists/INCOSE-REQUIREMENTS.md}).
El hallazgo más consistente: ningún requisito usa lenguaje normativo tipo ``shall'' (se usa formato
Historia de Usuario, una decisión de estilo consciente, no un descuido). La característica de conjunto
``verificable como conjunto'' ahora se cumple: los 8 de 8 requisitos Must tienen prueba automatizada real
(100\,\%, actualizado 2026-08-29 tras agregar prueba dedicada para RF-06 --- generación de actas, el
último que faltaba). A nivel de los 12 requisitos totales del sistema la cifra es 9/12 (75\,\%): RF-08,
RF-09 y RF-10 (prioridad Could/Should) siguen sin prueba dedicada --- no se declara 100\,\% de
trazabilidad general, solo de los Must que exige el criterio D0R.

\subsection{Trazabilidad}

\texttt{docs/trazabilidad/matriz.csv} conecta cada RF con su HU, módulo backend, endpoint REST,
prioridad MoSCoW, prueba automatizada real (o su ausencia declarada explícitamente) y evidencia
empírica. Esta matriz fue corregida en la Fase 7 del proyecto tras encontrarse que su versión anterior
citaba 5 archivos de prueba que nunca existieron y 3 referencias de evidencia empírica incorrectas (por
ejemplo, citar el puntaje de usabilidad general del sistema como evidencia de una funcionalidad CRUD
específica de gestión de salas) --- corrección verificable ejecutando
\texttt{scripts/validate-traceability.sh}.

\subsection{Bitácora de cambios y tasa de estabilidad}

\texttt{docs/requisitos/CHANGELOG-REQ.md} documenta la evolución de la SRS de la versión 0.9.0-rc (5 HU
formalizadas de 12 requisitos ya referenciados en la matriz) a la versión 1.0.0 (12 HU completas). La
tasa de estabilidad de \textit{alcance} es 1.00 (100\,\%): ningún requisito cambió su intención de
negocio entre versiones, solo se formalizaron 7 requisitos que ya existían operativamente. La tasa de
estabilidad de \textit{redacción} es 0.00 (0\,\%): se reescribió el 100\,\% del texto al nuevo formato.
Se reportan ambas cifras porque cualquiera de las dos por sí sola sería engañosa.

\subsection{Aprobación pendiente}

El SRS v1.0.0 incluye una página de aprobación lista para firma física o electrónica del
docente-director. A la fecha de este informe, esa firma \textbf{no se ha obtenido} --- es una acción que
requiere al docente-director, no generable automáticamente. Se declara explícitamente en vez de
simularse.
\newpage
